\chapter*{Abstract}
\addcontentsline{toc}{chapter}{Abstract}

The Linear Ordering Problem (LOP) is a classical NP-hard combinatorial optimisation
problem that consists of finding a permutation of rows and columns of a square matrix
such that the sum of superdiagonal entries is maximised. It has applications in various
fields including economics, archaeology, and scheduling.

In this work, we implement and evaluate iterative improvement (II) algorithms for the
LOP, combining two pivoting rules (first-improvement and best-improvement), three
neighbourhood structures (transpose, exchange, and insert), and two initialisation
methods (random permutation and the Chenery-Watanabe greedy heuristic), resulting in
twelve algorithm configurations. We further implement two Variable Neighbourhood
Descent (VND) algorithms that combine all three neighbourhoods in different orderings.

All neighbourhood evaluations use incremental delta computations to avoid redundant
objective function recalculations. Experiments are conducted on 78 benchmark instances
of sizes 150 and 250. Results are compared using the relative percentage deviation (RPD)
from best-known solutions and assessed through paired Wilcoxon signed-rank tests and
Student's t-tests.

The results show that insert-based and VND algorithms consistently outperform
transpose and exchange neighbourhoods. The initialisation method has a significant
impact for transpose and exchange, but not for insert. No statistically significant
difference was found between the two VND orderings.
